\newpage
\section{Societal Factors}

\subsection{Healthcare \& Rehabilitation}
In healthcare and rehabilitation, there is a wide need for orientation estimation from limb kinematics to assessment protocols for 
Axial Spondyloarthritis \fixme{add cites}. However, these studies
also highlight orientation drift due to IMUs as a limitation to their uses. 


Franco et al, highlights that IMUs have not widely been used for orientation estimation in a clincal setting due to gyroscopic drift \fixme{cite}.
Specifically, the study conducted range of motion tests that uses an IMU and Kalman Filter under constrained movements. The BASMI
movements performed and recorded all spanned for approximately 30s and achieved and RMSE = 2.0 degrees for joint flexion. 
However, all of these trials were performed
in a controlled laborartory setting with constrained movements to limit the accumulation of drift. Further, this
sensor alignment assumes that the sensors reflect the actual tilt of the
spine in the sagittal and frontal planes, and that the spine axial rotation is zero at t0 \fixme{cite}.
Under these constraints and assumptions, trying to extend this protocol outside of the laborartory and into clinical settings which have a lot 
of magnetic disturbances or under continuous monitoring of uncontrolled motions, their approach breaks down and drift accumulates. The proposed solution
in \fixme{section} serves as a way of extending and enabling further investigation of range of motion under non-constrained and continuous monitoring.
However, it is important to acknowledge that this proposed solution is based on data found in \fixme{section} compared to spinal kinematics data. However,
the dataset does include slow small rotations which is representative of the types of motions that Franco et al performed.  


\subsection{Agriculture}
\subsection{}

\begin{comment}
    Sy et al \fixme{cite} highlight that IMUs suffer from global drift 
and that these can be corrected with other sensors such as camera-based solutions but also mention the extra need of people needing to carry these 
cameras or that it requires setup \fixme{cite}. Wittmann et al also highlights the need for a homogenous magnetic field to correct gyroscope
drift \fixme{cite}. 
\end{comment}